\section{Introduction}
AlphaGo Zero and AlphaZero learn to play board games from nothing but the rules, by alternating self-play
with tree search and training a policy/value network on the resulting games
\citep{silver2017alphagozero,silver2018alphazero}.
The original systems used thousands of accelerators, but the algorithm is simple enough that small
versions are widely built: in courses, in hobby engines, and as research prototypes before a large run.
At that scale the budget for \emph{evaluation} is tiny, and the most common progress signal is also the
cheapest one: after each generation, play a short match between the new network and the previous one,
convert the score to an Elo difference \citep{elo1978rating}, and add it to a running total. We call this
total \emph{chained Elo}.

This paper reports a run where chained Elo told the wrong story. We trained a 9$\times$9 Go engine for
\valNGens{} generations on one laptop CPU. Chained Elo peaked at generation~\valChainedPeakGen{}
(\valChainedPeak{}) and then drifted down, putting generation~34 (\valChainedThirtyFour{}) below both
generation~25 (\valChainedTwentyFive{}) and generation~19 (\valChainedNineteen{}). Taken at face value,
training had stalled or regressed. In direct 80-game matches, however, generation~34 beat generation~25
in \valHTtfVsTf{} games and generation~19 in \valHTtfVsNt{}, and the final generation~\valNGens{} beat
generation~34 in \valHFfVsTtf{}.

\paragraph{Contributions.}
\begin{itemize}
  \item An evaluation of a complete small run against fixed anchors (\valHtoHGames{} games), showing steady
        improvement where chained Elo shows a decline (Figure~\ref{fig:elo}).
  \item An analysis of why chained Elo failed. Match noise grows with the number of generations, but the
        observed drop lies \valDropZ{} standard deviations from what noise around the true gain would produce. Replaying the same
        consecutive pairs at the anchored settings (64 rather than 32 simulations) gives gains that agree with
        the anchored matches (Figure~\ref{fig:step}). From this we derive a cheap protocol for small runs
        (Section~\ref{sec:lessons}).
  \item An open AlphaZero pipeline that trains on one laptop CPU. Its rules and search were checked to
        give identical results in Python, Swift with Core ML, and the browser with ONNX Runtime, so the
        same engine ships as a native app and a web page (Table~\ref{tab:platform}).
  \item An external calibration against GNU~Go, which shows how far such a network remains from a
        classical engine (Section~\ref{sec:gnugo}).
\end{itemize}
